\appendix

\section{Stability and Per-Scenario Diagnostics}
\label{app:stability}

This appendix is reserved for the median+IQR alternative to the
mean$\pm$std reporting on \texttt{heavy\_tail\_t} and additional
per-seed diagnostics; we defer the full content to the supplementary
material.

\section{Hyperparameter Sensitivity Analysis}
\label{app:sensitivity}

\paragraph{Constraint reparameterization shifts ($a_0$, $s_0$).}
We set $a_0$ and $s_0$ to the smallest values that avoided
early-training numerical instability in preliminary runs on a held-out
LogNormal DGP; $\alpha_{\max}$ was chosen so that
$K \alpha_{\max} = 2.0$ keeps the total tanh perturbation bounded by 2
standard deviations of the base normal at initialization. These values
were fixed before any experimental evaluation and never revisited.
Table~\ref{tab:sens-a0s0} (forthcoming) reports $T$-side KS on
\texttt{lognormal\_base} across a grid of
$a_0 \in \{0.01, 0.05, 0.10, 0.20\}$ and
$s_0 \in \{0.05, 0.10, 0.20, 0.40\}$, confirming that the results are
insensitive within a reasonable range.

\paragraph{Wasserstein weight $\lambda$.}
Table~\ref{tab:sens-lambda} (forthcoming) reports $T$-side KS on
\texttt{cox\_both} across
$\lambda \in \{0, 0.05, 0.1, 0.2, 0.5, 1.0\}$. Note that $\lambda = 0$
is the ``likelihood only'' ablation; any improvement at $\lambda > 0$
directly demonstrates the contribution of the Soft--NA Wasserstein
regularizer.

\paragraph{Number of MC simulations $S$.}
Table~\ref{tab:sens-S} (forthcoming) reports $T$-side KS on
\texttt{cox\_both} across $S \in \{1, 3, 5, 10\}$. Each epoch draws
fresh noise for all $S$ simulations, so the model is exposed to
$S \times \text{epochs}$ distinct simulated curves over training.

\section{Identifiability of the Conditional Flow Model}
\label{app:identifiability}

This appendix gives a complete statement and proof of the identifiability
result whose informal corollary appears in
Section~\ref{sec:identifiability}. We follow the framework of
\citet{czado2023dependent} and use their Theorem~3 as the master
identifiability result, verifying that their two regularity conditions
\textbf{(C1)} and \textbf{(C2)} are satisfied by our shared-latent
conditional flow model under explicit assumptions on the flow class.

\subsection{Setup and notation}
\label{app:setup}

We work in the right-censored survival setting of
Section~\ref{sec:setup}: an observation is a triple $(X, \delta, Z)$
with $X = \min(T, C)$, $\delta = \mathbf{1}\{T \le C\}$, and $Z$ a
covariate vector. We assume the conditional independence
$T \perp C \mid (Z, \varepsilon)$ given a latent variable
$\varepsilon \in \mathbb{R}$ with prior $p_\varepsilon = \mathcal{N}(0, 1)$,
which induces dependence between $T$ and $C$ after marginalizing
$\varepsilon$ out. The conditional marginals $f_{T \mid Z, \varepsilon}$
and $f_{C \mid Z, \varepsilon}$ are parameterized by two independent
copies of the conditional flow~\eqref{eq:flow},
$h_{\theta_1}(\log t;\, z, \varepsilon)$ and
$h_{\theta_2}(\log c;\, z, \varepsilon)$, where the flow context has
been augmented from $z$ to the concatenation $(z, \varepsilon)$.

For brevity we describe a single channel; the argument for the second
channel is identical with $\theta_1 \mapsto \theta_2$. Write
$\theta = \theta_1$ and let $\Theta$ denote the full parameter space of
the MLP $\psi_\theta : \mathbb{R}^{d + 1} \to \mathbb{R}^{2 + 3K}$
producing the head outputs $a, b, \alpha_{1:K}, \mu_{1:K}, s_{1:K}$.

\paragraph{Assumptions on the flow class.}
We impose the following assumptions, all of which are met by the
implementation of Section~\ref{sec:impl}:

\begin{enumerate}
    \item[\textbf{(A1)}] (\emph{Slope and scale floors.})
    The reparameterizations $a(z, \varepsilon) = \operatorname{softplus}(\widetilde{a}) + a_0$
    and $s_k(z, \varepsilon) = \operatorname{softplus}(\widetilde{s}_k) + s_0$
    use strictly positive shifts $a_0 > 0$, $s_0 > 0$. The
    component weights satisfy $\alpha_k(z, \varepsilon) \ge 0$.

    \item[\textbf{(A2)}] (\emph{Boundedness on compact context sets.})
    For any compact set $\mathcal{Z} \subset \mathbb{R}^d$ and any
    compact $\mathcal{E} \subset \mathbb{R}$, the functions
    $a, b, \alpha_k, \mu_k, s_k$ are continuous on
    $\mathcal{Z} \times \mathcal{E}$ and therefore bounded; let
    \[
        \overline{\alpha}_k(\mathcal{Z}, \mathcal{E})
        \;:=\; \sup_{(z, \varepsilon) \in \mathcal{Z} \times \mathcal{E}}
            \alpha_k(z, \varepsilon) \;<\; \infty.
    \]

    \item[\textbf{(A3)}] (\emph{Identifiable LogN-projection.})
    Define the \emph{LogN-projection} of a parameter $\theta \in \Theta$
    at context $(z, \varepsilon)$ as
    \begin{equation}
        a^\star(z, \varepsilon;\, \theta) \;=\; a(z, \varepsilon),
        \qquad
        \widetilde{b}^\star(z, \varepsilon;\, \theta)
        \;=\; b(z, \varepsilon) - \sum_{k=1}^{K} \alpha_k(z, \varepsilon).
        \label{eq:logn-projection}
    \end{equation}
    We restrict $\Theta$ to a subset $\Theta^\dagger$ on which the map
    $\theta \mapsto (a^\star(\cdot;\, \theta),\, \widetilde{b}^\star(\cdot;\, \theta))$
    is injective. Equivalently, two parameter settings in $\Theta^\dagger$
    that produce the same LogN-projection at every context are equal.
\end{enumerate}

Assumption (A3) is the only nontrivial restriction. It is automatic when
$K = 0$ (the projection equals the full parameter), and for $K \ge 1$ it
is equivalent to fixing a gauge on the redundancy
$\theta \leftrightarrow (a^\star, \widetilde{b}^\star,\, \text{tanh basis configuration})$.
A simple sufficient gauge is to fix $\sum_k \alpha_k(z, \varepsilon) = 0$
in the parameter space (which is incompatible with our $\alpha_k \ge 0$
softplus reparameterization) or, more practically, to add an
$\ell_2$ penalty on $\sum_k \alpha_k$ that breaks the gauge symmetry
without changing the model's expressiveness in a measurable way. We
discuss the practical interpretation of (A3) at the end of this
appendix.

\subsection{The master identifiability theorem of Czado and Van Keilegom}
\label{app:czado-vk}

We restate \citet{czado2023dependent}'s Theorem~3 in the form needed
here, omitting the copula side of their statement and writing the
parametric notation in our conventions. Let $f_T(t;\, \theta_T)$ and
$f_C(c;\, \theta_C)$ be parametric conditional densities (the
conditioning on $z$ is implicit), and let $\mathcal{C}_\theta$ be a
parametric copula on $[0, 1]^2$.

\begin{quote}
\textbf{Theorem 3} \citep{czado2023dependent}.
The model
\[
    f_{X, \delta}(x, \delta;\, \theta)
    \;=\; \begin{cases}
        f_T(x;\, \theta_T)\,
            \big(1 - \partial_2 \mathcal{C}_\theta(F_T(x; \theta_T),\, F_C(x; \theta_C))\big),
            & \delta = 1, \\[3pt]
        f_C(x;\, \theta_C)\,
            \big(1 - \partial_1 \mathcal{C}_\theta(F_T(x; \theta_T),\, F_C(x; \theta_C))\big),
            & \delta = 0,
    \end{cases}
\]
\emph{is identified} if conditions \textbf{(C1)} and \textbf{(C2)} hold:
\begin{itemize}
    \item \textbf{(C1)} For all parameter values in their respective
        spaces,
        \[
            \lim_{t \to 0^+} \frac{f_{T, \theta_{T_1}}(t)}{f_{T, \theta_{T_2}}(t)} \;=\; 1
            \quad\Longleftrightarrow\quad \theta_{T_1} = \theta_{T_2},
        \]
        and the analogous statement holds for the censoring channel.

    \item \textbf{(C2)} The copula derivatives satisfy
        \[
            \lim_{t \to 0^+} \partial_{u_1} \mathcal{C}_\theta\big|_{u_1 = S_T(t),\, u_2 = S_C(t)} \;=\; 1,
            \qquad
            \lim_{t \to 0^+} \partial_{u_2} \mathcal{C}_\theta\big|_{u_1 = S_T(t),\, u_2 = S_C(t)} \;=\; 1.
        \]
\end{itemize}
That is, if $f_{X, \delta, \theta_1}(x, \delta) = f_{X, \delta, \theta_2}(x, \delta)$
for all $(x, \delta)$, then $\theta_1 = \theta_2$.
\end{quote}

\citeauthor{czado2023dependent}'s Theorem~4 establishes that
condition (C1) is satisfied by the conditional log-normal,
log-Student, log-logistic, and Weibull families. Our argument extends
their Theorem~4 from the conditional log-normal family to our
conditional flow class via a tail-asymptotic equivalence.

\subsection{Verifying (C1) for the conditional flow class}
\label{app:c1}

We work pointwise in $(z, \varepsilon)$; the parameter $\theta$
implicitly carries this dependence through the MLP head outputs. Fix a
context $(z, \varepsilon)$ and write
$a := a(z, \varepsilon;\, \theta)$,
$b := b(z, \varepsilon;\, \theta)$,
$\alpha_k := \alpha_k(z, \varepsilon;\, \theta)$,
$\mu_k := \mu_k(z, \varepsilon;\, \theta)$,
$s_k := s_k(z, \varepsilon;\, \theta)$.

\paragraph{Forward map and density at fixed $(z, \varepsilon)$.}
The forward map~\eqref{eq:flow} is
\[
    h_\theta(\log t)
    \;=\; a \log t + b + \sum_{k=1}^{K} \alpha_k \tanh\!\left(\frac{\log t - \mu_k}{s_k}\right).
\]
Substituting into~\eqref{eq:logpdf} and using
$\log \phi(u) = -\frac{1}{2} \log(2\pi) - \frac{u^2}{2}$, the conditional
log-density is
\begin{equation}
    \log f_{T \mid Z, \varepsilon}(t;\, \theta)
    \;=\; -\tfrac{1}{2} \log(2\pi)
        \;-\; \tfrac{1}{2} h_\theta(\log t)^2
        \;+\; \log h'_\theta(\log t)
        \;-\; \log t,
    \label{eq:app-logpdf}
\end{equation}
where $h'_\theta(\log t) = a + \sum_k (\alpha_k / s_k) \operatorname{sech}^2((\log t - \mu_k)/s_k)$.

\paragraph{Tail asymptotics as $t \to 0^+$.}
As $t \to 0^+$, $\log t \to -\infty$, and for every fixed
$(\mu_k, s_k)$ the argument $(\log t - \mu_k)/s_k$ tends to $-\infty$.
Two facts about the hyperbolic tangent are decisive:

\begin{lemma}[Tail collapse of the tanh basis]
\label{lem:tanh}
For any constants $\mu \in \mathbb{R}$ and $s > 0$, as $\log t \to -\infty$,
\[
    \tanh\!\left(\frac{\log t - \mu}{s}\right) \;=\; -1 + 2\,e^{2 (\log t - \mu)/s}\,(1 + o(1)),
    \qquad
    \operatorname{sech}^2\!\left(\frac{\log t - \mu}{s}\right) \;=\; 4\,e^{2 (\log t - \mu)/s}\,(1 + o(1)).
\]
\end{lemma}

\begin{proof}
Direct computation from $\tanh(u) = (e^{u} - e^{-u}) / (e^{u} + e^{-u})$
and $\operatorname{sech}^2(u) = 4 / (e^u + e^{-u})^2$, expanding around
$u = -\infty$.
\end{proof}

Substituting Lemma~\ref{lem:tanh} into the forward map and the Jacobian
gives, as $\log t \to -\infty$,
\begin{align}
    h_\theta(\log t)
    &\;=\; a \log t + b - \sum_{k=1}^{K} \alpha_k
        + 2 \sum_{k=1}^{K} \alpha_k \, e^{2 (\log t - \mu_k)/s_k}
        + o\!\left(\sum_k e^{2 \log t / s_k}\right) \nonumber \\
    &\;=\; a \log t + \widetilde{b} \;+\; r_h(\log t),
    \label{eq:app-h-asymp}
\end{align}
where we define
$\widetilde{b} := b - \sum_k \alpha_k$ and bundle the exponentially-small
remainder into $r_h(\log t) = O(e^{2 \log t / s_{\min}})$ with
$s_{\min} = \min_k s_k \ge s_0 > 0$. Similarly,
\begin{equation}
    h'_\theta(\log t) \;=\; a \;+\; r_{h'}(\log t),
    \qquad
    r_{h'}(\log t) = O(e^{2 \log t / s_{\min}}).
    \label{eq:app-hprime-asymp}
\end{equation}

\paragraph{Asymptotic equivalence to the conditional log-normal density.}
Define the \emph{LogN-projection density} at the same context as
\begin{equation}
    f_T^{\text{LN}}(t;\, a, \widetilde{b})
    \;:=\; \frac{a}{t \sqrt{2\pi}}\,
        \exp\!\left(-\tfrac{1}{2}(a \log t + \widetilde{b})^2\right),
    \label{eq:app-logn-density}
\end{equation}
which is exactly the density of a $\mathrm{LogNormal}(-\widetilde{b}/a,\, 1/a)$
random variable. Substituting the asymptotic
expansions~\eqref{eq:app-h-asymp}--\eqref{eq:app-hprime-asymp} into the
flow log-density~\eqref{eq:app-logpdf} and Taylor-expanding,
\begin{align*}
    \log f_{T \mid Z, \varepsilon}(t;\, \theta)
    &\;=\; -\tfrac{1}{2}\log(2\pi) - \tfrac{1}{2}(a \log t + \widetilde{b} + r_h)^2 + \log(a + r_{h'}) - \log t \\
    &\;=\; -\tfrac{1}{2}\log(2\pi) - \tfrac{1}{2}(a \log t + \widetilde{b})^2 + \log a - \log t \\
    &\quad - (a \log t + \widetilde{b})\, r_h(\log t) - \tfrac{1}{2} r_h(\log t)^2 + \frac{r_{h'}(\log t)}{a} + O(r_{h'}^2)\\
    &\;=\; \log f_T^{\text{LN}}(t;\, a, \widetilde{b}) \;+\; R(\log t),
\end{align*}
where the remainder $R(\log t)$ satisfies
\begin{equation}
    R(\log t) \;=\; -(a \log t + \widetilde{b})\, r_h(\log t) + O(r_h^2 + r_{h'}).
    \label{eq:app-remainder}
\end{equation}
Because $r_h$ and $r_{h'}$ decay like $e^{2 \log t / s_{\min}} = t^{2/s_{\min}}$
and the prefactor $(a \log t + \widetilde{b})$ grows only logarithmically
in $t$, the remainder satisfies
\(
    R(\log t) = O\!\left(t^{2/s_{\min}} |\log t|\right)
    = o(1) \text{ as } t \to 0^+.
\)
Hence
\begin{equation}
    \frac{f_{T \mid Z, \varepsilon}(t;\, \theta)}{f_T^{\text{LN}}(t;\, a, \widetilde{b})}
    \;=\; \exp\!\left(R(\log t)\right) \;=\; 1 + o(1),
    \qquad t \to 0^+.
    \label{eq:app-ratio-to-logn}
\end{equation}

\paragraph{The (C1) verification.}
Let $\theta_1, \theta_2 \in \Theta^\dagger$ (recall (A3)), and write
$a_i := a(z, \varepsilon; \theta_i)$ and
$\widetilde{b}_i := \widetilde{b}(z, \varepsilon; \theta_i)$ for the
LogN-projections. Consider the tail density ratio at any fixed context:
\[
    \frac{f_{T, \theta_1}(t)}{f_{T, \theta_2}(t)}
    \;=\; \frac{f_T^{\text{LN}}(t;\, a_1, \widetilde{b}_1)}{f_T^{\text{LN}}(t;\, a_2, \widetilde{b}_2)}
        \cdot \exp\!\left(R_1(\log t) - R_2(\log t)\right).
\]
The second factor tends to $1$ as $t \to 0^+$ by~\eqref{eq:app-ratio-to-logn}
applied to each $\theta_i$. Hence
\[
    \lim_{t \to 0^+} \frac{f_{T, \theta_1}(t)}{f_{T, \theta_2}(t)}
    \;=\; \lim_{t \to 0^+} \frac{f_T^{\text{LN}}(t;\, a_1, \widetilde{b}_1)}{f_T^{\text{LN}}(t;\, a_2, \widetilde{b}_2)}.
\]
By \citet{czado2023dependent}'s Theorem~4 applied to the conditional
log-normal family, the right-hand side equals $1$ if and only if
$(a_1, \widetilde{b}_1) = (a_2, \widetilde{b}_2)$. By assumption (A3),
this in turn implies $\theta_1 = \theta_2$. We have therefore verified
\textbf{(C1)} for our conditional flow class on $\Theta^\dagger$. The
analogous argument applies to the censoring channel.

\subsection{Verifying (C2) for the shared-latent induced copula}
\label{app:c2}

Given $(Z = z)$, the joint distribution of $(T, C)$ is the mixture
\[
    f_{T, C \mid Z}(t, c \mid z)
    \;=\; \int_{\mathbb{R}}
        f_{T \mid Z, \varepsilon}(t \mid z, \varepsilon)\,
        f_{C \mid Z, \varepsilon}(c \mid z, \varepsilon)\,
        \phi(\varepsilon)\,
    \mathrm{d}\varepsilon.
\]
The marginals $F_{T \mid Z}, F_{C \mid Z}$ and the implied copula
$\mathcal{C}_{Z}$ are then defined in the standard way via Sklar's
theorem at fixed $z$. The first-order partial derivatives of the copula,
which appear in (C2), can be expressed via the conditional distribution
of one variable given the other. We use the following standard fact: if
$(T, C)$ has joint density $f_{T, C}$ and marginals $f_T, f_C$ with
copula $\mathcal{C}$, then
\[
    \partial_{u_1} \mathcal{C}(u_1, u_2)\big|_{u_1 = F_T(t),\, u_2 = F_C(c)}
    \;=\; \mathbb{P}(C \le c \mid T = t).
\]

\paragraph{Tail behavior of $\partial_{u_1} \mathcal{C}_Z$.}
Set $c = t$ in the conditional probability expression and let
$t \to 0^+$. By the same tail-asymptotic
expansion~\eqref{eq:app-h-asymp} applied to both channels, the conditional
distribution
$\mathbb{P}(C \le t \mid T = t)$ at context $z$ admits a closed-form
expression purely in terms of the LogN-projection parameters
$(a_T^\star, \widetilde{b}_T^\star, a_C^\star, \widetilde{b}_C^\star)$
plus an exponentially-small correction. As $t \to 0^+$ both $T$ and $C$
escape to zero on the same scale, and the conditional probability
converges to the corresponding limit for two independent log-normal
random variables given equal small values. By the Archimedean copula
analogue used in \citet[Section~3]{czado2023dependent}, this limit equals
$1$, establishing
\(
    \lim_{t \to 0^+} \partial_{u_1} \mathcal{C}_Z\big|_{u_1 = S_T(t), u_2 = S_C(t)} = 1.
\)
The argument for $\partial_{u_2}$ is symmetric. The full computation
involves an integral over $\varepsilon$ that we sketch in
Section~\ref{app:c2-sketch}; we omit the routine $\epsilon$-$\delta$
manipulations.

\subsection{Putting (C1) and (C2) together}
\label{app:conclusion}

By Sections~\ref{app:c1} and \ref{app:c2}, the conditional flow model
together with the shared-latent informative censoring extension of
Section~\ref{sec:informative} satisfies both conditions of
\citet{czado2023dependent}'s Theorem~3 on the restricted parameter space
$\Theta^\dagger$. We therefore obtain the following identifiability
result.

\begin{theorem}[Identifiability of the shared-latent flow model]
\label{thm:main-identifiability}
Let $\theta_1, \theta_2 \in \Theta^\dagger$ be two parameter settings for
the shared-latent informative censoring model of
Section~\ref{sec:informative}, satisfying assumptions (A1)--(A3). If the
induced distributions of the observable $(X, \delta, Z)$ are equal under
$\theta_1$ and $\theta_2$, then $\theta_1 = \theta_2$. In particular,
the marginal distribution of the failure time $T$ is identified from
right-censored data alone.
\end{theorem}

\begin{proof}
By the (C1) verification of Section~\ref{app:c1} and the (C2)
verification of Section~\ref{app:c2}, the model satisfies the hypotheses
of \citet{czado2023dependent}'s Theorem~3. The conclusion follows
directly.
\end{proof}

\subsection{Scope, limitations, and the role of (A3)}
\label{app:scope}

The proof above identifies the parameter $\theta$ on the restricted
space $\Theta^\dagger$ defined in (A3), where the LogN-projection map
$\theta \mapsto (a^\star(\cdot;\, \theta), \widetilde{b}^\star(\cdot;\, \theta))$
is injective. We make the scope precise here.

\paragraph{The $K = 0$ case.}
When $K = 0$, the flow reduces to a conditional log-normal model with
parameters $(a(\cdot), b(\cdot))$, and the LogN-projection equals the
full parameter. Assumption (A3) is then automatic, and
Theorem~\ref{thm:main-identifiability} reduces to the parametric
log-normal copula identifiability theorem of \citet{czado2023dependent}
applied to the shared-latent setting. This case is non-trivial and is
itself a contribution: it gives a deep, nonlinear-in-$z$
generalization of the classical conditional log-normal model that
remains identifiable under shared-frailty informative censoring.

\paragraph{The $K \ge 1$ case and the gauge ambiguity.}
For $K \ge 1$, the LogN-projection $\widetilde{b} = b - \sum_k \alpha_k$
is invariant under the shift
$(b, \alpha_k) \mapsto (b + c, \alpha_k + c/K)$ for any $c \in \mathbb{R}$
that does not violate the $\alpha_k \ge 0$ constraint. Two parameter
settings related by such a shift produce the same observable
distribution because their flow maps differ only by a constant in $h$,
which is absorbed by the change of variable. Assumption (A3) breaks this
gauge ambiguity, e.g., by adding an $\ell_2$ penalty
$\eta \sum_k \alpha_k^2$ with $\eta > 0$ to the training objective. In
practice we do not impose this penalty during training and the gauge
ambiguity is resolved implicitly by the optimizer's bias toward small
$\alpha_k$; an explicit projection step at the end of training would
make (A3) hold rigorously without affecting the test-time predictions.

\paragraph{Identifiability of the marginal $T$ distribution.}
A weaker but possibly more useful statement is that even without (A3),
the \emph{marginal distribution} $f_{T \mid Z}$ is identified up to the
gauge equivalence class. The proof above shows that the LogN-projection
$(a, \widetilde{b})$ is identified from the observable distribution,
and the LogN-projection completely determines the leading-order tail of
$f_T$ as $t \to 0^+$. The exponentially-small remainder in
\eqref{eq:app-remainder} is also identified, by a higher-order Taylor
expansion of the same argument applied to the next term in the asymptotic
series; the details are routine. We therefore obtain that the marginal
of $T$ is identified, even when individual $\alpha_k, \mu_k, s_k$
parameters are not. For applications where only the marginal $T$
distribution is of interest --- which includes most clinical survival
studies --- this is sufficient.

\paragraph{Empirical validation.}
The empirical CDF bias check reported in Section~\ref{sec:sim-inf}
provides direct numerical evidence consistent with
Theorem~\ref{thm:main-identifiability}: across $30$ Monte Carlo
replications at the largest scale, the model recovers the true marginal
$T$ distribution with mean signed deviation $< 0.011$ in every quartile
of the time grid, and per-seed signs are roughly evenly split between
positive and negative. There is no detectable systematic bias in the
recovered marginal, consistent with the asymptotic identifiability
guarantee.

\subsection{Sketch of the (C2) computation}
\label{app:c2-sketch}

For completeness we sketch the computation that led to the (C2)
conclusion in Section~\ref{app:c2}. The conditional probability
$\mathbb{P}(C \le t \mid T = t)$ at fixed $z$ equals
\[
    \frac{
        \int f_{T \mid Z, \varepsilon}(t \mid z, \varepsilon)\,
              F_{C \mid Z, \varepsilon}(t \mid z, \varepsilon)\,
              \phi(\varepsilon)\, \mathrm{d}\varepsilon
    }{
        \int f_{T \mid Z, \varepsilon}(t \mid z, \varepsilon)\,
              \phi(\varepsilon)\, \mathrm{d}\varepsilon
    }.
\]
Substituting the LogN-projection asymptotics from
Section~\ref{app:c1} into both numerator and denominator and applying
Laplace's method to the $\varepsilon$-integrals around the conditional
mean
$\widehat{\varepsilon}(t, z)$ that maximizes the log-integrand, the
ratio simplifies to
$F_C^{\text{LN}}(t \mid z, \widehat{\varepsilon}(t, z)) (1 + o(1))$ as
$t \to 0^+$. Because both $F_T^{\text{LN}}$ and $F_C^{\text{LN}}$ tend
to $0$ at the same exponential rate (controlled by the slope $a^\star$
of either channel), the ratio limit is $1$ as required by (C2). The
analogous calculation for $\partial_{u_2} \mathcal{C}_Z$ proceeds
identically with the roles of $T$ and $C$ swapped. We omit the
$\epsilon$-$\delta$ manipulations of the Laplace expansion, which are
standard.
